\section{The normalized input to the filling theorem}\label{cs:filling}
The hypotheses of \cite[Theorem 1.1]{Oscillatory} now hold for the
actual sequence \eqref{sa:coefficient}, with
\begin{equation}\label{sa:export}
 \varepsilon_N=(-1)^{dR},\qquad
 \alpha=2R-(R-1)/2=(3R+1)/2,\qquad
 a_N(\delta)=\frac{(2\pi)^{-(R-1)/2}}{\sqrt{\det H_N(\delta)}}.
\end{equation}
The limiting $a_N$ is a nonzero holomorphic function. The finite
functions $A_N$ must remain inside the exponential; their convergence
does not justify replacing $e^{NA_N}$ by $e^{NA_0}$ in a relative formula.

For fixed independently signed sufficiently large powers $n_e$, that
theorem applies to the original central Fourier period with scalar
$\Gamma_N=\exp(\pi\ii d^2\sum n_e/(4N))$. The product of its parity
factor and \eqref{sa:export} is
$(-1)^{(N+1)R+dR}=1$ at every integer level. If $z_N$ is the
unit-charge spectral saddle, $S_N$ its finite action, and $Q_N$ its
negative spectral Hessian, the resulting leading expression is
\[
 \Gamma_N(-2\ii)^R\sqrt{2\pi}\,N^{R+1/2}
 \frac{\prod_e\sin(2\pi z_{N,e})}
       {\sqrt{\det H_N(z_N)}\sqrt{\det Q_N}}e^{NS_N}.
\]
The relative error and eventual nonvanishing are precisely the
conclusions of the oscillatory theorem. Its proof supplies the whole
charge sum and endpoint currents; none is inferred from the present
face Gaussian alone. Exterior framing or component units, and the
modified-dimension factor when that convention is used, remain those
of \cite{FinitePaper}.

A geometric companion may identify $2\pi\Phi$ with an augmented volume
and $2\pi\Re S_0$ with a specified filled volume. Those statements
require its marked peripheral and metric proofs. A full invariant
limit also requires control of the complementary degree sectors.
Neither conclusion is part of Theorem~\ref{sa:main}.

